\subsection{Breadboard}
\markright{Breadboard}

\subsubsection*{Definition}
A breadboard (or solderless breadboard) is a reusable construction base for prototyping
electronic circuits without soldering \cite{wiki:breadboard}. Beneath the plastic
surface, rows of spring-contact strips connect groups of holes electrically, allowing
component leads and jumper wires to be inserted and removed at will.

\labimage{lab-01/assets/breadboard.jpg}%
  {Full-size solderless breadboards (830 tie-points) with power rails
   (top and bottom) and component strips in the center section.}%
  {breadboard}

\subsubsection*{Specifications}
Holes are spaced at 0.1\,inch (2.54\,mm) pitch to match DIP integrated circuit packages.
A standard full-size breadboard provides 830 tie-points. Contact resistance is a few
ohms; inter-strip parasitic capacitance is a few picofarads, limiting reliable operation
to below roughly 10\,MHz \cite{analog:breadboard}.

\subsubsection*{Purpose of Use}
The breadboard enables rapid assembly, testing, and modification of circuits during the
design and verification phase. No tools are required: components are inserted by hand,
and reconfiguring a circuit is a matter of moving wires rather than removing solder.

\subsubsection*{Applications}
Breadboards are the standard prototyping platform in electronics education and in early
product development. They are used for assembling analog filters, digital logic circuits,
microcontroller (Arduino, Raspberry Pi) projects, and sensor interface circuits
\cite{analog:breadboard}. Once validated, designs move to veroboard or custom PCB.
